(1) 高斯光束的复参数在 ABCD 系统下满足 \(q_{\rm out}=(Aq_{\rm in}+B)/(Cq_{\rm in}+D)\)。薄透镜后
\[
q_0=\frac{iz_R}{1-iz_R/f}
=\frac{-z_R^2/f+iz_R}{1+(z_R/f)^2}.
\]
实部为负、虚部为正，表示束腰将在透镜后某正距离处形成。

(2) 自由传播 s 后 \(q(s)=q_0+s\)。束腰处复 q 的实部必须为0，因此
\[
s_w=-\Re q_0=\frac{fz_R^2}{f^2+z_R^2}.
\]
因 f,zR>0，该位置在透镜后方。

(3) 新 Rayleigh 长度为 \(z'_R=\Im q_0=z_R/[1+(z_R/f)^2]=f^2z_R/(f^2+z_R^2)\)。束腰半径满足 \(w_0'^2=\lambda z_R'/\pi\)，故
\[
\frac{w_0'}{w_0}=\sqrt{\frac{z'_R}{z_R}}
=\frac f{\sqrt{f^2+z_R^2}}.
\]
这两个量可从复参数的实、虚部同时核验。

(4) f≫zR 时 \(s_w\simeq z_R^2/f\) 很短，\(z_R'\simeq z_R,w_0'\simeq w_0\)，弱透镜基本不改变束腰。f≪zR 时 \(s_w\simeq f,\ z_R'\simeq f^2/z_R,\ w_0'/w_0\simeq f/z_R\)，强聚焦产生远小于入射束腰的新腰斑。
